% 目录中附录编号前统一加“附录”前缀
\addtocontents{toc}{%
  \protect\renewcommand{\protect\cftchappresnum}{\appendixname\protect\space}%
  \protect\renewcommand{\protect\cftchapaftersnum}{}%
}

\chapter{检查、烧录与常见问题}
\label{app:concise-check}

\section{检查与摘要}

向导每一步都会实时检查当前配置，进入“检查与摘要”页时集中列出全部问题。问题分为两类：

\begin{itemize}
  \item \textbf{错误}：配置不完整或自相矛盾，阻止生成。常见原因包括通道号未填、同一
        引脚分配给两个设备、按键重复分配、必填项为空；
  \item \textbf{警告}：配置合法但可能不符合预期，只作提醒。常见原因包括冲刺速度偏低、
        拨弹电机位于低频引脚组、摩擦轮最大转速偏低、速度与引脚角色不完全匹配。
\end{itemize}

每条问题都带有定位信息，点击可直接跳回对应页面修改。修改后返回检查页，计数随之更新。
“配置概览”汇总通道、速度、引脚用途等关键配置，供生成前最后核对。

警告不等于错误，但也不应一律忽略。出现警告时，应确认当前方案确实接受该结果，再继续
生成；无法判断时，按警告提示调整配置。

\section{编译与烧录}

“烧录”页负责编译与写入。编译器一栏保持默认的内置编译器即可；已安装其他编译器的
使用者可以在下拉框中切换，但首次使用不建议改动。

页面提供三个操作：

\begin{itemize}
  \item \textbf{仅编译}：只完成编译，不写入主控板，用于确认配置与编译环境正常；
  \item \textbf{编译并烧录}：完成编译并写入已连接的主控板；
  \item \textbf{导出固件}：把编译结果保存为文件，供其他工具使用。
\end{itemize}

烧录前关闭主控板 SERVO、POWER 与机械拓展板 POWER、BOOSTER 四个供电开关，将主控板
断电后重新通过 USB 连接。页面检测到主控板后即可烧录；烧录成功后设备自动复位，检测到的
烧录设备消失属于正常现象。再次烧录前需要重新断电、上电。

Linux 首次烧录需要按页面提示安装设备访问权限；安装完成后重新插拔主控板。烧录失败时，
页面下方的进度与日志保留失败原因，可对照常见问题排查。

\begin{note}
只能烧录主控板。机械拓展板不运行用户固件，向拓展板烧录会破坏其出厂程序。烧录期间
不要插拔接线，也不要切换供电开关。
\end{note}

\section{装车前检查单}

以下检查单适用于步兵、工程与调试项目。每一项都能直接观察或操作，全部通过后才落地运行。

\begin{enumerate}
  \item 车体架空，轮子离地；弹仓清空，射界内无人；
  \item 关闭主控板 SERVO、POWER 与机械拓展板 POWER、BOOSTER 四个供电开关后连接 USB；
  \item 只向主控板烧录，确认烧录日志显示成功；
  \item 上电后听到启动提示音，确认程序已运行；
  \item 逐轮检查转向与整车方位一致，异常时修改配置而非现场改接线；
  \item 逐轴检查云台舵机方向与行程，确认不触碰机械限位；
  \item 检查拨弹电机方向与单发时长，确认卡弹时可反向退弹；
  \item 检查摩擦轮启动、加速与停止是否平滑，确认两侧转速无明显差异；
  \item 让测试序列或动作循环两轮，确认无偶发异常；
  \item 断电后再改接线或换设备；落地后先低速、空载运行，再逐步提速。
\end{enumerate}

\begin{note}
安全底线：不向机械拓展板烧录；不在上电状态插拔接线；摩擦轮最大转速不超过 800；
机构出现异响、堵转或超出行程时立即切断电源。机械限位是最后一道保护，任何时候都不能
省略。
\end{note}

\section{常见问题}

\begin{table}[H]
  \centering
  \small
  \begin{tabularx}{\linewidth}{@{}>{\raggedright\arraybackslash}p{0.28\linewidth}
      >{\raggedright\arraybackslash}X@{}}
    \toprule
    现象 & 处理 \\
    \midrule
    车完全不动 & 确认手柄与主控板均已上电；核对通道号是否一致；检查死区是否过大；确认
      烧录成功 \\
    前进却后退 & 摇杆极性可能与配置相反；架空低速核对，必要时调整方向或转向设置 \\
    个别轮子反向 & 修改该轮的“方向”，不改接线；四个轮子同时反向时先核对整机前后约定 \\
    转向相反 & 打开“反转左右转向”，前后移动不受影响 \\
    拨弹一卡一卡 & 拨弹电机位于低频引脚组，属正常现象；提高拨弹速度、缩短单发时长，
      或改用连续拨弹 \\
    摩擦轮带不动 & 检查最大转速是否偏低、机械阻力是否过大；电调首次使用需完成油门
      行程学习 \\
    未检测到主控板 & 关闭四个供电开关；主控板断电后重新上电；更换 USB 口与数据线，
      避免经过集线器 \\
    烧录中断或写块无响应 & 更换 USB 口与数据线；关闭占用设备的其他程序；重新断电
      上电后重试 \\
    首次编译失败 & 检查项目路径与缓存目录是否包含中文字符；把工具目录加入杀毒软件
      白名单后重试 \\
    Web 版不能烧录 & 浏览器无法访问本地设备；配置与检查可在 Web 版完成，编译烧录
      回到桌面版 \\
    模式切换后轮子停 & 新模式没有映射这四路输出；在动作映射中补全该模式需要的输出 \\
    旧项目打不开 & 旧版本存档可能不被当前版本接受；用对应旧版打开并另存，或按新版本
      重建配置 \\
    \bottomrule
  \end{tabularx}
  \caption{常见问题速查}
  \label{tab:concise-faq}
\end{table}
